% =====================================================================
% CSE 341: Concepts of Programming Languages
% Project D5 (Part 1): Retrospective & Self-Assessment for ANS
% Student: 220104004926 Yusuf Hamza Çelebi
% Compile with: tectonic 220104004926_D5_P1.tex
% =====================================================================
\documentclass[11pt,a4paper]{article}

\usepackage[utf8]{inputenc}
\usepackage[T1]{fontenc}
\usepackage{mathpazo}
\usepackage[scaled=0.92]{helvet}
\usepackage[varqu,varl]{zi4}
\usepackage[english]{babel}
\usepackage{microtype}
\usepackage[a4paper,margin=2.25cm,top=2.0cm,bottom=2.15cm]{geometry}

\usepackage[table]{xcolor}
\definecolor{ansblue}{HTML}{1E3A8A}
\definecolor{anscaption}{HTML}{6B7280}
\definecolor{anscodebg}{HTML}{F9FAFB}
\usepackage[colorlinks=true, linkcolor=ansblue, urlcolor=ansblue,
            citecolor=ansblue, pdfborder={0 0 0}]{hyperref}

\usepackage[explicit]{titlesec}
\titleformat{\section}
  {\large\bfseries\sffamily\color{black}}
  {\thesection}{0.65em}{#1}
\titlespacing*{\section}{0pt}{11pt}{5pt}
\setlength{\parskip}{4.5pt plus 1pt minus 1pt}
\setlength{\parindent}{0pt}
\linespread{1.08}
\setcounter{secnumdepth}{0}

\usepackage{fancyhdr}
\pagestyle{fancy}
\fancyhf{}
\renewcommand{\headrulewidth}{0pt}
\fancyfoot[L]{\footnotesize\color{anscaption}%
  CSE 341: ANS Retrospective (D5, Part 1) \(\cdot\)
  Student 220104004926 Yusuf Hamza Çelebi}
\fancyfoot[R]{\footnotesize\color{anscaption}Page \thepage}

\newcommand{\code}[1]{\texttt{#1}}

\hypersetup{
  pdftitle  = {ANS Retrospective and Self-Assessment (D5, Part 1)},
  pdfauthor = {220104004926 Yusuf Hamza Çelebi},
  pdfsubject= {CSE 341 Project Part 1},
  pdfkeywords={DSL, chemistry, retrospective, self-assessment, ANS}
}

\begin{document}

\begin{center}
  {\footnotesize\color{anscaption}CSE 341: Concepts of Programming Languages\par}
  \vspace{3pt}
  {\huge\bfseries\sffamily ANS: Alchemist's Notation Script\par}
  \vspace{4pt}
  {\large Retrospective \& Self-Assessment (D5), Part 1\par}
  \vspace{6pt}
  {\footnotesize\color{anscaption}
    Student: 220104004926 Yusuf Hamza Çelebi \(\cdot\)
    Spring 2026 \(\cdot\)
    Gebze Technical University\par}
\end{center}
\vspace{7pt}

\section{What Is Working}

The lexer and parser seem to be working. They passed the current tests, but we still need to be careful because more cases can appear later.

\section{What Is Hard}

Determining whether the lexical structure is consistent with the EBNF, and implementing that exact lexical structure into the lexer and parser, was difficult. C is one of the main system programming languages, so we decided to implement the lexer and parser in C.

\section{Plan For Part 2}

We are planning to fix the output statement. We first decided to implement output as a reserved keyword. But we think the \code{print} keyword is kind of irrelevant to the overall syntax of ANS. We will change \code{print} to a function-style construct like \code{mol()}.

\end{document}
